\documentclass[11pt]{article}
\usepackage[utf8]{inputenc}
\usepackage{amsmath,amssymb,amsthm}
\usepackage[margin=1in]{geometry}
\usepackage{microtype}
\usepackage{hyperref}
\usepackage{xcolor}
\usepackage{enumitem}
\usepackage{newunicodechar}
\emergencystretch=3em
\newunicodechar{ℝ}{\ensuremath{\mathbb{R}}}
\newunicodechar{λ}{\ensuremath{\lambda}}
\newunicodechar{κ}{\ensuremath{\kappa}}
\newunicodechar{·}{\ensuremath{\cdot}}
\newunicodechar{∂}{\ensuremath{\partial}}

\newcommand{\deriv}{\operatorname{deriv}}
\newcommand{\tideref}[2][]{\href{https://therisensea.org/tides/#2/}{\ifx&#1&\texttt{#2}\else#1\fi}}

\title{Tide retrospective:\\\emph{the fourth connected cumulant (flow-equation identity)}}
\author{Claude Opus 4.8 (1M ctx), with Daniel Murfet}
\date{2026-05-30}

\begin{document}
\maketitle

\begin{abstract}
The capstone of the anharmonic cumulant arc: the fourth connected cumulant of
$u=-(tx)$ as the third derivative of the perturbed mean at $h=0$,
\[
  \kappa_4 = \partial_h^3\big|_0 \langle u\rangle_h
    = \langle u^4\rangle_0 - 4\langle u^3\rangle_0\langle u\rangle_0
      - 3\langle u^2\rangle_0^2 + 12\langle u^2\rangle_0\langle u\rangle_0^2
      - 6\langle u\rangle_0^4,
\]
in closed $(G_0,\dots,G_4)$ form. Three derivative layers --- mean
$M=G_1/G_0$, susceptibility $S=\deriv M$, third cumulant $S_2=\deriv S$ ---
stitched with \texttt{Filter.EventuallyEq.deriv}. New file
\texttt{AnharmonicFourthCumulant.lean}; builds within the default heartbeat
budget, no \texttt{sorry}/\texttt{axiom}/\texttt{native\_decide}.
\end{abstract}

\section{Setting}

The closing tide of a single-day anharmonic arc on the laplace seabed, and the
flow-equation / fourth-cumulant identity flagged as the headline target of the
day's candidate survey. The cumulant ladder was built rung by rung: the mean
($K'$), the susceptibility $\kappa_2$ ($\partial_h$ of the mean), its general-$h$
version, and $\kappa_3$ ($\partial_h^2$ of the mean). $\kappa_4$ is
$\partial_h^3$ of the mean --- one differentiation layer beyond $\kappa_3$,
which is what made it the heaviest single tide.

\section{The thing formalised}

Writing $G\,n$ for \texttt{weightedPartition lam alpha gamma t n}:
\begin{quote}\itshape
For the anharmonic potential, $\lambda,\gamma>0$, $\alpha^2<3\lambda\gamma$,
$t>0$,
\begin{align*}
  &\mathrm{iteratedDeriv}\ 3\ (h\mapsto G\,1\,h/G\,0\,h)\ 0 \\
  &\qquad = \frac{G_4 G_0^3 - 4G_3 G_1 G_0^2 - 3G_2^2 G_0^2 + 12 G_2 G_1^2 G_0 - 6G_1^4}{G_0^4},
\end{align*}
all $G_n$ at $0$ --- the standard moments-to-cumulants combination after
dividing each $G_n(0)$ by $G_0(0)$.
\end{quote}
A companion theorem, \texttt{anharmonic\_susceptibility\_deriv\_general}, gives
the third cumulant at \emph{any} $|h_0|<1$ (the general-$h$ $\deriv S$), which
is what the chaining consumes.

\section{Proof strategy}

Three layers. The mean's derivative is the susceptibility $S$ (general-$h$
susceptibility tide); $S$'s derivative is $S_2=(G_3G_0^2-3G_1G_2G_0+2G_1^3)/G_0^3$
(proved here at general $h_0$ by a compound quotient rule); and $S_2$'s
derivative at $0$ is $\kappa_4$ (a second compound quotient rule, the numerator
built from \texttt{.mul}/\texttt{.pow}/\texttt{.const\_mul}/\texttt{.sub}/\texttt{.add}).
The three are joined: $\deriv M =_{\text{ev}} S$ and $\deriv S =_{\text{ev}} S_2$
on $\mathrm{ball}\,0\,1$, so by \texttt{Filter.EventuallyEq.deriv} the second
derivatives agree eventually too, and
$\mathrm{iteratedDeriv}\,3\,M\,0 = \deriv(\deriv(\deriv M))\,0 = \deriv S_2\,0
= \kappa_4$.

\section{Roadblocks and resolutions}

Two, both mechanical. \emph{(i) \texttt{set\_option} placement.} The
\texttt{maxHeartbeats} bumps I added pre-emptively sat between docstring and
theorem, which is a parse error (the docstring must abut the declaration); they
turned out to be unnecessary anyway --- the degree-eight \texttt{ring} closes
within the default budget --- so they were simply removed. \emph{(ii) Pi-vs-%
pointwise.} Building the numerator's \texttt{HasDerivAt} from sub-terms like
\texttt{hG0.pow 2} leaves \texttt{Pi.pow}/\texttt{Pi.mul} \emph{applications}
($\,(G_0^2)\,0\,$ rather than $\,(G_0\,0)^2\,$) inside the derivative value, which
\texttt{ring} treats as opaque atoms distinct from their pointwise forms. An
ascription of the top-level function fixes that function but not the value;
a \texttt{simp only [Pi.pow\_apply, Pi.mul\_apply]} before \texttt{field\_simp}
pushes the applications inward and lets \texttt{ring} finish. (The third-%
cumulant tide avoided this only because its sub-terms happened to stay
pointwise.)

\section{What was Mathlib, what was new}

\textbf{Mathlib.} The \texttt{HasDerivAt} combinators (\texttt{.mul},
\texttt{.sub}, \texttt{.add}, \texttt{.pow}, \texttt{.const\_mul}, \texttt{.div},
\texttt{.deriv}); \texttt{iteratedDeriv\_succ}/\texttt{\_one};
\texttt{Filter.EventuallyEq.deriv}; the \texttt{Pi.\_apply} simp lemmas;
\texttt{field\_simp}, \texttt{ring}.

\textbf{Seabed (reused).} The pointwise partition derivatives, the general-$h$
susceptibility, and the $G_0(0)>0$ positivity --- all from this session's arc.

\textbf{New.} The general-$h$ third cumulant and the fourth-cumulant theorem.
About 160 lines.

\section{Lessons}

\begin{itemize}[itemsep=3pt]
\item \textbf{Ascribe the \emph{sub-terms}, not just the outer function.} A
  function ascription on a compound \texttt{HasDerivAt} does not stop
  \texttt{Pi.pow}/\texttt{Pi.mul} from appearing in the inferred derivative
  \emph{value}; a \texttt{simp only [Pi.\_apply]} before \texttt{ring} is the
  reliable cleanup for nested quotient/product rules.
\item \textbf{Try the default heartbeat budget before bumping it.} Even the
  degree-eight cumulant \texttt{ring} fit the default; the pre-emptive
  \texttt{maxHeartbeats} was dead weight (and triggers a style linter).
\item \textbf{The cumulant ladder is uniform.} $\kappa_2,\kappa_3,\kappa_4$ are
  the same recipe at increasing order: differentiate the previous cumulant
  function (needing its general-$h$ form) and chain. $\kappa_5$ would be one
  more identical layer.
\end{itemize}

\section{Follow-ups}

\begin{itemize}[itemsep=3pt]
\item \textbf{Connect to $\log Z$ explicitly.} Package $\kappa_2,\dots,\kappa_4$
  as $\partial_h^n \log Z|_0$ via the cumulant generating function, closing the
  loop with \texttt{logPartition\_hasDerivAt}.
\item \textbf{Moment-form corollaries.} State each $\kappa_n$ in the
  $\langle u^k\rangle_0$ (\texttt{gibbsExp}) form via moment-normalisation.
\item \textbf{Promote \texttt{amgm\_t\_abs\_x} to \texttt{lean-common}.}
\end{itemize}

\end{document}
